% Generated draft: requires mathematical and source review.
\documentclass[11pt]{article}
\usepackage[margin=25mm]{geometry}
\usepackage{fontspec}
\setmainfont{Latin Modern Roman}
\usepackage{amsmath,amssymb,mathtools,unicode-math}
\setmathfont{Latin Modern Math}
\usepackage{enumitem,needspace,xurl,hyperref}
\setlength{\parindent}{0pt}
\setlength{\parskip}{5pt}
\setlength{\emergencystretch}{3em}
\newcommand{\N}{\mathbb{N}}
\newcommand{\Z}{\mathbb{Z}}
\newcommand{\Q}{\mathbb{Q}}
\newcommand{\R}{\mathbb{R}}
\newcommand{\C}{\mathbb{C}}
\newcommand{\F}{\mathbb{F}}
\newcommand{\A}{\mathbb{A}}
\newcommand{\PP}{\mathbb{P}}
\newcommand{\E}{\mathbb{E}}
\newcommand{\eps}{\varepsilon}
\newcommand{\dd}{\,\mathrm{d}}
\begin{document}
\section*{UnsolvedMath ID: 701641}
\textbf{Source problem code:} catalogue entry
\subsection{Definitions and mathematical statement}
Let \(\N=\{1,2,3,\ldots\}\). For a prime \(p\), let \(C_p\) denote the cyclic group of order \(p\), and let
\[
W_p=C_p\wr C_p=(C_p)^p\rtimes C_p
\]
be the regular wreath product, in which the final \(C_p\) acts by cyclically permuting the \(p\) coordinates. An epimorphic image of a group means the image of a surjective group homomorphism; thus \(G\) has no epimorphic image isomorphic to \(W_p\) precisely when no normal subgroup \(N\trianglelefteq G\) satisfies \(G/N\cong W_p\).

For this draft, interpret the \emph{\(p\)-lower central series} as the lower exponent-\(p\) central series
\[
P_1(G)=G,\qquad P_{n+1}(G)=P_n(G)^p[P_n(G),G]\quad(n\geq1).
\]
Here \(H^p\) is the subgroup generated by the \(p\)-th powers of elements of \(H\), and \([H,G]\) is generated by the commutators \([h,g]=h^{-1}g^{-1}hg\), with \(h\in H\) and \(g\in G\). The supplied record does not define its series convention; this interpretation requires confirmation against the original source.

For a set \(S\) of primes, consider the assertion
\[
\begin{gathered}
\exists f:\N\times\N\longrightarrow\N\quad
\forall p\in S\quad\forall d\in\N\quad
\forall\text{ finite }p\text{-groups }G:\\
\left( G\text{ is generated by }d\text{ elements and }
\nexists N\trianglelefteq G\text{ with }G/N\cong W_p\right)
\Longrightarrow
\left(\forall n\geq1,\quad
|P_n(G)/P_{n+1}(G)|\leq f(p,d)\right).
\end{gathered}
\]
Thus the bound may depend on \(p\) and \(d\), but must be independent of \(G\) and the series index \(n\). The source asks for this assertion in the two cases
\[
\text{(a) }S=\{p:p\text{ is prime and }p\geq7\},
\qquad
\text{(b) }S=\{2,3,5\}.
\]
The phrase ``with \(d\) generators'' is interpreted here as admitting a generating set of size at most \(d\); the record does not explicitly require \(d\) to be the minimal number of generators.

\textbf{Supplied status qualification.} The registry labels the problem partially solved. Its supplied discussion reports that part (a) is refuted for \(p\geq7\), attributing an AI solution to M. Kinyon and J. Urban and reporting confirmation by P. Spiga. This is a supplied source assessment, not an independently verified mathematical conclusion. Part (b) is retained as the remaining conjectural question in that record. The source also invites comparison with a characterization by Aner Shalev of finitely generated \(p\)-adic analytic pro-\(p\) groups, without supplying the characterization or a precise citation.
\subsection{Short English statement}
P. Spiga asks whether excluding the quotient \(C_p\wr C_p\) forces every factor of the \(p\)-lower central series of a finite \(d\)-generated \(p\)-group to have order bounded solely in terms of \(p\) and \(d\). The supplied record reports a refutation for primes \(p\geq7\) and retains the cases \(p=2,3,5\) as conjectural. The precise series convention requires source confirmation.
\subsection{Sources}
\begin{itemize}
\item Supplied local UnsolvedMath registry record, permanent ID \(701641\), exact source code \texttt{catalogue entry}, titled \emph{Kourovka Notebook Problem 19.93}. Proposed by P. Spiga in 2018. The supplied registry status is \texttt{partially\_solved}.
\item \emph{The Kourovka Notebook (Unsolved Problems in Group Theory)}, 19th issue (2018), Problem 19.93; also the supplied 21st issue edition (2026), identified as \texttt{arXiv:1401.0300v48}. The collection description names E. I. Khukhro and V. D. Mazurov as editors of Issues 1--21 (1965--2026).
\item M. Kinyon and J. Urban, \emph{Preprint} (2026), supplied link: \url{https://kourovkanotebookorg.wordpress.com/wp-content/uploads/2026/09/kourovka-mathematics.pdf}. The supplied discussion reports that part (a) is refuted for \(p\geq7\), describes the result as an AI solution, and reports confirmation by P. Spiga. No preprint title is supplied.
\item Supplied status-review date: \texttt{2026-10-16}; evidence classification: \texttt{maintained\_source}. The review attributes the solved marking for part (a) to the October 2026 21st issue, \texttt{arXiv:1401.0300v48}. This assessment is supplied dated curation and has not been independently verified here.
\item The supplied statement mentions Aner Shalev's characterization of finitely generated \(p\)-adic analytic pro-\(p\) groups as background for comparison; no precise bibliographic reference or theorem statement is supplied.
\end{itemize}
\end{document}
